\documentclass[10pt,a4paper]{amsart}
\usepackage[margin=19mm]{geometry}
\usepackage{amsmath,amssymb,mathtools}
\usepackage[scr=rsfs,cal=euler]{mathalfa}
\usepackage{xfrac}
\usepackage[unicode,hidelinks,bookmarks=false]{hyperref}
\newcommand{\ie}{i.e.\ }
\newcommand{\cf}{cf.\ }
\newcommand{\resp}{respectively\ }
\newcommand{\Subsection}{\subsection}
\providecommand{\nobreakdash}{\nobreak\hbox{-}}
\setlength{\emergencystretch}{2em}
\multlinegap0pt
\usepackage{xcolor}
\usepackage{amssymb}
\usepackage{mathtools}
\usepackage{bm}
\usepackage[shortlabels]{enumitem}
\usepackage{tikz}
\usepackage{cancel}
\newtheorem{proposition}{Proposition}
\title{Positive Solutions for an Indefinite-Weight Minkowski Mean Curvature Neumann Problem: Multiplicity and Asymptotic Behaviour}
\author{Ricardo Ziegele}
\date{Source: arXiv:2609.09894v1 (CC BY 4.0)}
\begin{document}
\maketitle

\noindent\emph{Source and licence.} Positive Solutions for an Indefinite-Weight Minkowski Mean Curvature Neumann Problem: Multiplicity and Asymptotic Behaviour. Version arXiv:2609.09894v1 (CC BY 4.0), distributed by arXiv under the Creative Commons Attribution 4.0 International licence, http://creativecommons.org/licenses/by/4.0/, as recorded in the arXiv licence metadata for this exact version and shown on the abstract page https://arxiv.org/abs/2609.09894v1. Full licence notice: \texttt{licenses/arxiv-2609.09894-CC-BY-4.0.txt}.\par\smallskip

\noindent\textit{Editorial scope.} Counted unit: the article's proposition prop:inactive-region (source0.tex lines 1268--1296). The article's complete original proof of it is delivered with it (source0.tex lines 1298--1627, one \texttt{proof} environment of 330 lines). No mathematics is written, repaired or re-derived: the statement and the proof are the article's own lines, in the article's own notation, and no retained line is substituted or re-flowed.  Retained uncounted as the source-local context the proof needs: lines 1195--1195; lines 1197--1198.  Nothing was omitted from any retained span, so the package carries no omission record.  No retained line contains a citation, so the wrapper carries no bibliography and no bibliography entry.  No retained line contains a figure environment, an image-inclusion command or drawing code, so no image is delivered and the package declares no asset dependency.  Editorial formatting only, in addition: an AMS \texttt{amsart} preamble that restates the article's own declarations and macros used by the retained lines, namely the environments \texttt{proposition} on separate counters, together with the standard packages those lines need.  The retained environments print in this document as Proposition 1 in order of appearance, whereas the article prints the same items with the numbering of its own full document.  The complete native source of the article as distributed in the arXiv e-print for this exact version is bundled unchanged as \texttt{sources/209858/source0.tex}.
\begin{equation} \label{eq:defJC} \mathcal{J}(u) := \int_\Omega a(x) |u|^{p} \, dx, \quad \mathcal{C}(u) := \int_\Omega a(x) |u|^{p-2}u \, dx,\end{equation}

\begin{equation}\label{eq:P} \tag{P}  \max_{u \in W^{1,\infty}(\Omega)} \mathcal J(u) \quad \text{s.t.} \,\quad \mathcal{C}(u) = 0, \,   \|\nabla u \|_{L^{\infty}(\Omega)} \leq 1  .
\end{equation}

\begin{proposition}
\label{prop:inactive-region}
Let $u \in K_a$ be a maximizer of \eqref{eq:P}. Let
$U\subset\Omega$ be a connected open set such that
$$
\|\nabla u\|_{L^\infty(U)}<1,
$$
and
$$
a(x)\neq 0
\qquad\text{for a.e. }x\in U.
$$
Then $u$ is constant in $U$. More precisely, there exists
$\mu\in\mathbb R$ independent of $U$, such that
$$u\equiv0
\qquad\text{or}\qquad
u\equiv\frac{p-1}{p}\mu
\quad\text{in }U.$$
Moreover, whenever $u\not\equiv0$ in $U$, the constant $\mu$ is
characterized by
$$
D\mathcal J(u)[\varphi]
=
\mu D\mathcal C(u)[\varphi]
\qquad
\text{for every }\varphi\in C_c^\infty(U),
$$
where $\mathcal J$ and $\mathcal C$ are as in \eqref{eq:defJC}.
\end{proposition}

\begin{proof}
To begin with, if $u\equiv0$ in $U$, then there is nothing to prove.
We therefore assume that $u\not\equiv0$ in $U$. Since $p>2$, the map
$s\mapsto |s|^{p-2}s$ belongs to $C^1(\mathbb R)$, and hence
$\mathcal C$ is of class $C^1$ on $W^{1,\infty}(\Omega)$. Its Fr\'echet
derivative is given by
$$
D\mathcal C(u)[\psi]
=
(p-1)\int_\Omega
a(x)|u|^{p-2}\psi\,dx.
$$
Similarly, for $\mathcal{J}$, we have
$$
D\mathcal J(u)[\psi]
=
p\int_\Omega
a(x)|u|^{p-2}u\,\psi\,dx.
$$
Since $u\not\equiv0$ in $U$ and $a\neq0$ a.e. in $U$, we have
$$
a(x)|u(x)|^{p-2}\not\equiv0
\qquad\text{in }U.
$$
Consequently, there exists
$$
\psi\in C_c^\infty(U)
$$
such that
\begin{equation}\label{eq:Cnotnull}
D\mathcal C(u)[\psi]\neq0.
\end{equation}
Indeed, otherwise
$$
\int_U a(x)|u|^{p-2}\psi\,dx=0
\qquad
\text{for every }\psi\in C_c^\infty(U),
$$
and this would imply
$$
a(x)|u(x)|^{p-2}=0
\qquad\text{a.e. in }U,
$$
a contradiction. Now fix an arbitrary $\varphi\in C_c^\infty(U)$ and consider
$$
F(t,s)
:=
\mathcal C(u+t\varphi+s\psi).
$$
Since
$$
F(0,0)=\mathcal C(u)=0
$$
and
$$
\partial_sF(0,0)
=
D\mathcal C(u)[\psi]\neq0,
$$
the implicit function theorem yields the existence of a $\varepsilon>0$ and a
$C^1$ function
$$
s:(-\varepsilon,\varepsilon)\to\mathbb R,
\qquad
s(0)=0,
$$
such that
$$
\mathcal C
\bigl(
u+t\varphi+s(t)\psi
\bigr)=0
$$
for every $|t|<\varepsilon$. Moreover, differentiating this identity at $t=0$ gives
$$
D\mathcal C(u)[\varphi]
+
s'(0)D\mathcal C(u)[\psi]
=0,
$$
and hence
\begin{equation} \label{eq:sprime}
s'(0)
=
-
\frac{
D\mathcal C(u)[\varphi]
}{
D\mathcal C(u)[\psi]
}.
\end{equation}
We claim that, after possibly reducing $\varepsilon$, the curve
$$
u_t:=u+t\varphi+s(t)\psi \in K_a \quad \text{for every } |t| < \varepsilon. 
$$
Indeed, since
$$
\left(\operatorname{supp}\varphi\cup\operatorname{supp}\psi \right)
\subset U
$$
and $\|\nabla u\|_{L^\infty(U)}<1$, there exists $\delta>0$ such that
$$
\|\nabla u\|_{L^\infty(U)}
\le 1-\delta.
$$
Since $s(0)=0$ and $s\in C^1$, we have $s(t)=O(t)$ as $t\to0$. Therefore, for $|t|$ sufficiently small,
$$
|s(t)|\leq C|t|
$$
for some $C>0$. Choosing $|t|$ so small that
$$
|t|\bigl(\|\nabla\varphi\|_{L^\infty(\Omega)}+C\|\nabla\psi\|_{L^\infty(\Omega)}\bigr)
\leq \frac{\delta}{2},
$$
we obtain
$$
\|\nabla u_t\|_{L^\infty(U)}
\leq 1-\frac{\delta}{2}<1.
$$
On the other hand, since $\operatorname{supp}\varphi\cup\operatorname{supp}\psi\subset U$, one has $u_t=u$ on $\Omega\setminus U$, and hence
$$
|\nabla u_t|=|\nabla u|\leq1
\qquad\text{a.e. in }\Omega\setminus U.
$$
Consequently $\|\nabla u_t\|_{L^\infty(\Omega)}\leq1$. Together with $\mathcal C(u_t)=0$, this proves that $u_t\in K_a$. Since $u$ is a maximizer of $\mathcal J$ on $K_a$, the function $t\longmapsto\mathcal J(u_t)$ has a local maximum at $t=0$. Hence
$$
0
=
\left.\frac{d}{dt}\right|_{t=0}
\mathcal J(u_t)
=
D\mathcal J(u)
\bigl[
\varphi+s'(0)\psi
\bigr].
$$
Using \eqref{eq:sprime}, we obtain
$$
D\mathcal J(u)[\varphi]
=
\mu
D\mathcal C(u)[\varphi] 
$$
where
\begin{equation} \label{eq:defmu}
\mu:=
\frac{
D\mathcal J(u)[\psi]
}{
D\mathcal C(u)[\psi]
} = \frac{p\int_\Omega a(x) |u|^{p-2}u \, \psi \,dx}{(p-1) \int_\Omega a(x) |u|^{p-2} \, \psi \, dx}.
\end{equation}
Since $\varphi\in C_c^\infty(U)$ was arbitrary,
$$
D\mathcal J(u)[\varphi]
=
\mu D\mathcal C(u)[\varphi]
\qquad
\forall\varphi\in C_c^\infty(U).
$$
Substituting the expressions for the derivatives yields
$$
p\int_U
a(x)|u|^{p-2}u\varphi\,dx
=
\mu(p-1)
\int_U
a(x)|u|^{p-2}\varphi\,dx
$$
for every $\varphi\in C_c^\infty(U)$. Therefore,
$$
a(x)|u|^{p-2}
\left(
pu-\mu(p-1)
\right)
=0
\qquad\text{a.e. in }U.
$$
Since $a\neq0$ a.e. in $U$, we conclude that
$$
|u|^{p-2}
\left(
pu-\mu(p-1)
\right)
=0
\qquad\text{a.e. in }U.
$$
Thus
$$
u(x)
\in
\left\{
0,\frac{p-1}{p}\mu
\right\}
\qquad\text{for a.e. }x\in U.
$$
Finally, due to $u\in W^{1,\infty}(\Omega)$, it admits a continuous representative. Then, considering that $U$ is connected, a continuous function which takes values only in the two-point set $\left\{
0,\frac{p-1}{p}\mu
\right\}
$ must be constant. Consequently,
$$
u\equiv0
\qquad\text{in }U,
$$
or
$$
u\equiv\frac{p-1}{p}\mu
\qquad\text{in }U,
$$
thus finishing this part of the proof. It remains to show that $\mu$ does not depend on $U$. Let $U_1, U_2$ be two open, connected subsets where  $u \not \equiv 0$ and $\|\nabla u\|_{L^\infty (U_i)} < 1$ for $i = 1, 2$, and let $\mu_1$ and $\mu_2$ be the respective constants given by the procedure above, applied in $U_1 $ and in $U_2$, respectively. If $U_1 \cap U_2 \not = \emptyset $, then, taking any $\varphi \in C_c^{\infty}(U_1 \cap U_2)$, one has
$$D \mathcal J (u) [\varphi] = \mu_1 D \mathcal{C}(u)[\varphi] \quad \text{and} \quad  D \mathcal J (u) [\varphi] = \mu_2 D \mathcal{C}(u)[\varphi]$$
which, by subtracting, yields
$$0 = (\mu_1 - \mu_2) D \mathcal{C}(u)[\varphi] \quad \forall \varphi \in C_c^{\infty}(U_1 \cap U_2),$$
and, since reasoning as before, we have that $D\mathcal C(u) \not \equiv 0$, it implies $\mu_1 = \mu_2$. So, let's assume that $U_1$ and $U_2$ are now disjoint, with associated
constants $\mu_1$ and $\mu_2$ and test functions
$\psi_1\in C_c^\infty(U_1)$ and $\psi_2\in C_c^\infty(U_2)$ as above,
satisfying \eqref{eq:Cnotnull}. Fix
$\varphi_2\in C_c^\infty(U_2)$, and this time define
$$
u_t:=u+t\varphi_2+s(t)\psi_1.
$$
The argument follows exactly as before, by applying the implicit function
theorem to the map
$$
F(t,s):=
\mathcal C(u+t\varphi_2+s\psi_1).
$$
Indeed,
$$
F(0,0)=\mathcal C(u)=0
$$
and
$$
\partial_sF(0,0)
=
D\mathcal C(u)[\psi_1]\neq0.
$$
Therefore, by the implicit function theorem, there exist
$\varepsilon>0$ and a $C^1$ function
$$
s:(-\varepsilon,\varepsilon)\to\mathbb R,
\qquad s(0)=0,
$$
such that
$$
\mathcal C(u+t\varphi_2+s(t)\psi_1)=0
$$
for every $|t|<\varepsilon$.

Moreover, since $U_1$ and $U_2$ are inactive regions and the perturbation is supported in $U_1\cup U_2$, the same local argument used above shows, after possibly reducing $\varepsilon$, that
$$
\|\nabla u_t\|_{L^\infty(U_1\cup U_2)}<1,
\qquad
|\nabla u_t|=|\nabla u|\leq1
\quad\text{a.e. in }\Omega\setminus(U_1\cup U_2).
$$
Hence $\|\nabla u_t\|_{L^\infty(\Omega)}\leq1$, and therefore $u_t\in K_a$ for every $|t|<\varepsilon$. Since $u$ is a maximizer of $\mathcal J$ on $K_a$, the function
$t\mapsto\mathcal J(u_t)$ has a local maximum at $t=0$. Consequently,
\begin{equation} \label{eq:4xx}
0
=
\left.\frac{d}{dt}\right|_{t=0}
\mathcal J(u_t)
=
D\mathcal J(u)[\varphi_2]
+s'(0)D\mathcal J(u)[\psi_1].
\end{equation}
On the other hand, differentiating
$$
\mathcal C(u+t\varphi_2+s(t)\psi_1)=0
$$
at $t=0$ gives
$$
D\mathcal C(u)[\varphi_2]
+s'(0)D\mathcal C(u)[\psi_1]
=0,
$$
and therefore
$$
s'(0)
=
-\frac{D\mathcal C(u)[\varphi_2]}
        {D\mathcal C(u)[\psi_1]}.
$$
Substituting this expression into \eqref{eq:4xx}, we obtain
$$
D\mathcal J(u)[\varphi_2]
=
\frac{D\mathcal J(u)[\psi_1]}
     {D\mathcal C(u)[\psi_1]}
D\mathcal C(u)[\varphi_2].
$$
By the definition of $\mu_1$,
$$
\mu_1
=
\frac{D\mathcal J(u)[\psi_1]}
     {D\mathcal C(u)[\psi_1]},
$$
and hence
$$
D\mathcal J(u)[\varphi_2]
=
\mu_1D\mathcal C(u)[\varphi_2]
\qquad
\forall\varphi_2\in C_c^\infty(U_2).
$$
However, by the definition of the multiplier $\mu_2$ associated with
$U_2$, we also have
$$
D\mathcal J(u)[\varphi_2]
=
\mu_2D\mathcal C(u)[\varphi_2]
\qquad
\forall\varphi_2\in C_c^\infty(U_2).
$$
Consequently,
\begin{equation} \label{eq:mu1mu2}
(\mu_1-\mu_2)
D\mathcal C(u)[\varphi_2]=0
\qquad
\forall\varphi_2\in C_c^\infty(U_2).
\end{equation}
But, once again,  by the definition of $\psi_2$ (cf. \eqref{eq:Cnotnull}), we know that
$$
D\mathcal C(u) \not \equiv 0,
$$
and so, \eqref{eq:mu1mu2} implies $\mu_1=\mu_2,$
finishing the proof.
\end{proof}

\end{document}
